\section{Discussion and Conclusion}
\label{sec:discussion}

We define indirect prompt injection by authority rather than wording: planted
content is an attack when it attempts to influence a decision its source is not
authorized to determine, even when it only tampers with the task in progress. Our
benchmark makes this boundary explicit in 90 validated
task-tampering cases across three environments; within each case, the planted
content justifies a fixed harmful outcome through different influence
primitives, so defenses can be compared per primitive, not only by aggregate
ASR. Our results show that defenses tested mainly against redirecting
instructions leave task tampering only partly covered: objective attacks (\(b(U)\)), the kind
prior evaluations mainly test, already succeed less often without a defense
than attacks on decision rules, available actions, and state, and \(b(U)\)
remains among the three lowest-ASR primitives under every defense but
CaMeL.

The results suggest two directions for defense design, both grounded in
authority. Content screening can miss attacks whose content reads as ordinary
task data, such as \(b(O)\) and \(b(s)\) attacks; screening could therefore also
ask whether a source is authorized to determine the targeted decision, not
only whether its content reads as a redirecting instruction. Tool-level gates
can fail to block attacked calls that use tools the task needs; action checks
could instead admit each consequential argument only from a source authorized
to set it.
